\chapter{Trích dẫn nguồn và từ chối câu hỏi ngoài corpus}
\label{ch:trich-dan}

\section{Vì sao chọn}

Câu trả lời về quy chế chỉ có giá trị khi người đọc kiểm tra lại được. Và với câu hỏi mà quy chế không
nói tới, ``không tìm thấy'' là câu trả lời đúng --- bịa ra một con số là lỗi nặng nhất.

\section{Cơ chế}

\paragraph{Trích dẫn.} \texttt{format\_sources()} đọc metadata của các Node nguồn qua \texttt{node\_label()}
và in ``Nguồn: <văn bản> -- Điều $N$''. Nguồn lấy từ metadata chứ không để LLM tự ghi, nên luôn khớp với
những gì đã vào prompt. Prompt còn bắt LLM trích nguyên văn câu quy định kèm tên văn bản và số Điều
trước khi kết luận.

\paragraph{Từ chối.} Hai lớp: (1) prompt cho LLM một lối ra hợp lệ --- ``nếu trích đoạn không có thông
tin, hãy nói không tìm thấy''; (2) \texttt{SimilarityPostprocessor} loại Node dưới ngưỡng
(\texttt{similarity\_cutoff}), đã cài nhưng mặc định tắt vì điểm cosine thấp và sát nhau
(Mục~\ref{sec:retriever}).

\section{Đã làm gì}

\begin{itemize}
  \item Nhãn nguồn thống nhất với văn bản embedding qua trường \texttt{nguon\_trich} (Chương~\ref{ch:node-parser}).
  \item Bộ \texttt{eval/cau\_hoi\_danh\_gia.json}: 13 câu, mỗi câu có \texttt{pham\_vi} (10 ``trong'', 3 ``ngoài'')
        để tách từ chối đúng khỏi từ chối sai.
\end{itemize}

\section{Số liệu}

Trên 56 câu hỏi tự đặt (Chương~\ref{ch:trien-khai}), tỷ lệ trả lời ``không tìm thấy'' là 52\% ở phiên
đầu và 41\% ở phiên chạy lại ngày 08/10. Chấm tay cho thấy cả hai loại: từ chối đúng (``điểm rèn luyện
là gì'', ``bao nhiêu tiền một tín chỉ'') và từ chối sai (``còn kỹ sư'', ``rớt một môn có được xuất sắc
không'').

\chuatrienkhai{Tỷ lệ từ chối đúng trên một bộ câu ngoài corpus đủ lớn (ít nhất 10 câu) chưa đo; bộ hiện
có mới 3 câu ``ngoài''. Trích dẫn chưa tới mức khoản/trang.}

\section{Dùng ở đâu trong đồ án}

Mọi câu trả lời ở CLI, giao diện Gradio và nhật ký đánh giá đều kèm dòng ``Nguồn''.
